% Hanzo Threshold Signing: MPC Key Management for AI Infrastructure
\documentclass[11pt]{article}
\usepackage{shared/hanzocover}
\usepackage{amsmath, amssymb, amsthm}
\usepackage{mathtools}
\usepackage{bm}
\usepackage{graphicx}
\usepackage{booktabs}
\usepackage{hyperref}
\usepackage{enumitem}
\usepackage{algorithm}
\usepackage{algpseudocode}
\usepackage{xcolor}
\usepackage{float}
\usepackage{multirow}

\hypersetup{colorlinks=true,linkcolor=black,citecolor=blue,urlcolor=blue}

\newtheorem{theorem}{Theorem}
\newtheorem{lemma}[theorem]{Lemma}
\newtheorem{proposition}[theorem]{Proposition}
\newtheorem{corollary}[theorem]{Corollary}
\newtheorem{definition}{Definition}
\newtheorem{remark}{Remark}

\title{HANZO THRESHOLD SIGNING\\[0.3em]
{\large MPC Key Management for\\AI Infrastructure}}
\author{
    Zach Kelling \\
    \textit{Hanzo Industries Inc (Techstars '17)} \\
    \texttt{research@hanzo.ai}
}
\date{March 2026}

\begin{document}
\hanzocoverpage

% ============================================================================
% ABSTRACT
% ============================================================================
\begin{abstract}
We present Hanzo Threshold Signing, a system for distributed key management and
cryptographic signing across Hanzo infrastructure services. The system eliminates
single points of compromise for critical signing operations: IAM session tokens,
KMS secret encryption keys, agent credential delegation, and ACI compute
attestations. We implement two threshold signature protocols: (i)~FROST
(Flexible Round-Optimized Schnorr Threshold signatures) for Ed25519, used in
IAM JWT signing and agent credential chains, and (ii)~CGGMP21 for threshold
ECDSA (secp256k1), used in blockchain transaction signing for ACI staking and
Lux cross-chain operations. The system operates as a $(t, n)$ threshold scheme
where $t$ of $n$ signers must participate to produce a valid signature, with no
single party ever holding the complete private key. We describe the distributed
key generation (DKG) protocol, the signing protocol with identifiable abort,
proactive key refresh without changing the public key, and integration with
Hanzo KMS for share storage. Production deployment across 5 signing nodes
achieves $(3, 5)$ threshold signing with median latency of 12ms for FROST and
45ms for CGGMP21, processing over 2 million threshold signatures per day.
\end{abstract}

% ============================================================================
\section{Introduction}
\label{sec:intro}
% ============================================================================

Hanzo infrastructure depends on cryptographic signing for security-critical
operations:
\begin{enumerate}[leftmargin=1.4em]
    \item \textbf{IAM session tokens}: JWT access tokens signed with Ed25519,
    issued by Hanzo IAM~\cite{hanzoiam} and validated by every service in the
    platform.
    \item \textbf{KMS envelope encryption}: Master keys that encrypt all
    secrets in Hanzo KMS~\cite{hanzokms}. Compromise of a master key exposes
    every secret in the system.
    \item \textbf{Agent credentials}: AI agents~\cite{hanzoagentsdk} receive
    delegated credentials for tool access, API calls, and resource management.
    The signing key authorizing these delegations must not be compromised.
    \item \textbf{ACI compute attestations}: GPU nodes attest to inference
    results by signing output hashes. These attestations are verified on-chain
    for the Hanzo ACI compute market~\cite{hanzoaci}.
\end{enumerate}

In all cases, a single signing key stored on a single server is a single point
of failure. If that server is compromised, the attacker can forge IAM tokens,
decrypt all secrets, issue fraudulent agent credentials, or submit false compute
attestations.

Threshold signatures solve this by distributing the signing key across $n$
servers such that any $t$ must cooperate to sign, but fewer than $t$ servers
learn nothing about the key. The combined public key is a standard Ed25519 or
ECDSA public key---verifiers need no knowledge of the threshold scheme.

\subsection{Protocol Selection}

We use two protocols, chosen for specific properties:

\begin{table}[H]
\centering
\small
\begin{tabular}{lll}
\toprule
\textbf{Protocol} & \textbf{Curve} & \textbf{Use Case} \\
\midrule
FROST~\cite{frost} & Ed25519 & IAM tokens, agent credentials, API signing \\
CGGMP21~\cite{cggmp21} & secp256k1 & Blockchain transactions, ACI staking \\
\bottomrule
\end{tabular}
\caption{Protocol selection by use case.}
\label{tab:protocols}
\end{table}

FROST is preferred for Ed25519 operations due to its simplicity (2-round
signing, no zero-knowledge proofs during signing). CGGMP21 is required for
ECDSA because Schnorr-based threshold constructions do not produce valid ECDSA
signatures.

% ============================================================================
\section{FROST for Ed25519}
\label{sec:frost}
% ============================================================================

\subsection{Distributed Key Generation}

FROST DKG produces key shares without any party learning the full secret key.
Each participant $P_i$ acts as both dealer and verifier in a Pedersen
DKG~\cite{pedersen1991}:

\begin{definition}[$(t, n)$-Secret Sharing]
A $(t, n)$-Shamir secret sharing of $s \in \mathbb{Z}_q$ distributes shares
$s_i = f(i)$ where $f(X) = s + \sum_{j=1}^{t-1} a_j X^j$ is a random
polynomial of degree $t-1$. Any $t$ shares determine $s$; fewer than $t$
shares reveal no information about $s$.
\end{definition}

\begin{algorithm}
\caption{FROST Distributed Key Generation}
\label{alg:dkg}
\begin{algorithmic}[1]
\Require Participants $P_1, \ldots, P_n$, threshold $t$
\Ensure Public key $Y$, shares $s_1, \ldots, s_n$
\For{each participant $P_i$}
    \State Sample random polynomial $f_i(X) = a_{i,0} + \sum_{j=1}^{t-1} a_{i,j} X^j$
    \State Broadcast commitments $C_{i,j} = a_{i,j} \cdot G$ for $j = 0, \ldots, t-1$
    \State Send share $f_i(k)$ privately to each $P_k$
\EndFor
\For{each participant $P_k$}
    \State Verify: $f_i(k) \cdot G \stackrel{?}{=} \sum_{j=0}^{t-1} C_{i,j} \cdot k^j$ for all $i$
    \State Compute combined share: $s_k = \sum_{i=1}^{n} f_i(k)$
\EndFor
\State $Y = \sum_{i=1}^{n} C_{i,0}$ \Comment{Combined public key}
\end{algorithmic}
\end{algorithm}

\begin{theorem}[DKG Security]
The FROST DKG protocol is secure against $t-1$ corrupted participants under
the discrete logarithm assumption on the Ed25519 curve. No coalition of fewer
than $t$ participants can reconstruct the secret key $s = \sum_i a_{i,0}$.
\end{theorem}

\subsection{Threshold Signing}

FROST signing requires two rounds. In the first round, each signer generates
and broadcasts a nonce commitment. In the second round, each signer produces a
partial signature using their key share and the aggregated nonce.

\begin{definition}[FROST Signature Share]
Given message $m$, group nonce $R = \sum_{i \in S} R_i$, and challenge
$c = H(R \| Y \| m)$, signer $P_i$ with share $s_i$ computes:
\begin{equation}
z_i = d_i + e_i \cdot \rho_i + \lambda_i \cdot s_i \cdot c,
\end{equation}
where $(d_i, e_i)$ are nonce secrets, $\rho_i = H(i, m, B)$ is the binding
factor, and $\lambda_i$ is the Lagrange interpolation coefficient for the
signing set $S$.
\end{definition}

The aggregated signature $\sigma = (R, z)$ where $z = \sum_{i \in S} z_i$ is a
standard Ed25519 signature verifiable with the public key $Y$.

\begin{proposition}[Unforgeability]
FROST signatures are existentially unforgeable under chosen-message attack
(EUF-CMA) in the random oracle model, assuming the hardness of the discrete
logarithm problem on Curve25519~\cite{frost}.
\end{proposition}

\subsection{Nonce Pre-computation}

To reduce signing latency, signers pre-compute nonce pairs during idle periods:
\begin{enumerate}[leftmargin=1.4em]
    \item Each signer generates a batch of $B$ nonce pairs $(d_j, e_j)$ and
    commitments $(D_j, E_j)$.
    \item Commitments are broadcast and stored by all participants.
    \item During signing, a pre-computed nonce pair is consumed. No
    communication is needed in the first round.
\end{enumerate}

With $B = 10000$ pre-computed nonces, signing reduces to a single round with
median latency of 8ms (vs.\ 12ms for 2-round signing).

% ============================================================================
\section{CGGMP21 for ECDSA}
\label{sec:cggmp}
% ============================================================================

ECDSA threshold signing is more complex than Schnorr because the ECDSA
equation $s = k^{-1}(H(m) + r \cdot x)$ involves multiplication of secret
values ($k^{-1}$ and $x$), requiring MPC multiplication protocols.

\subsection{Key Generation}

CGGMP21~\cite{cggmp21} key generation extends Feldman VSS with Paillier
encryption for MPC multiplication:
\begin{enumerate}[leftmargin=1.4em]
    \item Each party generates a Paillier key pair $(N_i, \phi_i)$ with
    $N_i > 2^{2048}$ and proves correctness via a zero-knowledge proof of
    safe prime product.
    \item Shamir sharing of the ECDSA private key $x$ proceeds as in FROST DKG.
    \item Each party encrypts their share under their Paillier key and
    broadcasts the ciphertext with a range proof.
\end{enumerate}

\subsection{Threshold Signing Protocol}

The CGGMP21 signing protocol consists of four phases:

\begin{enumerate}[leftmargin=1.4em]
    \item \textbf{Nonce generation}: Each signer $P_i$ samples $k_i, \gamma_i$
    and broadcasts commitments $\Gamma_i = \gamma_i \cdot G$ and Paillier
    encryptions $K_i = \texttt{Enc}(k_i)$.
    \item \textbf{MPC multiplication}: Parties compute $\delta_i = k_i \gamma_j
    + k_j \gamma_i$ pairwise using the Paillier homomorphism, with
    affine-group range proofs ensuring honest behavior.
    \item \textbf{Nonce reveal}: Parties reveal $\delta_i$ and reconstruct
    $\delta = k \cdot \gamma$. Compute $R = (\delta^{-1}) \cdot \Gamma$ where
    $\Gamma = \sum_i \Gamma_i$. Extract $r = R.x \bmod q$.
    \item \textbf{Partial signatures}: Each party computes $\sigma_i = k_i
    \cdot H(m) + k_i \cdot r \cdot \chi_i$ where $\chi_i = x_i \cdot
    \lambda_i$ is the Lagrange-weighted share.
\end{enumerate}

\begin{theorem}[Identifiable Abort]
If any party deviates from the CGGMP21 protocol, the honest parties can
identify the cheater with probability 1 and exclude them from future signing
sessions. This relies on the zero-knowledge range proofs attached to every
Paillier operation.
\end{theorem}

% ============================================================================
\section{Proactive Key Refresh}
\label{sec:refresh}
% ============================================================================

Key shares must be periodically refreshed to limit the window of vulnerability
if a share is compromised. Proactive secret sharing~\cite{herzberg1995}
re-randomizes shares without changing the public key.

\begin{definition}[Key Refresh]
A key refresh transforms shares $(s_1, \ldots, s_n)$ of secret $s$ into fresh
shares $(s_1', \ldots, s_n')$ of the same secret $s$, such that knowing $s_i$
and $s_j'$ (for different epochs) provides no advantage in recovering $s$.
\end{definition}

The protocol:
\begin{enumerate}[leftmargin=1.4em]
    \item Each party $P_i$ generates a random polynomial $g_i(X)$ of degree
    $t-1$ with $g_i(0) = 0$.
    \item $P_i$ sends $g_i(k)$ privately to each $P_k$.
    \item Each party computes $s_k' = s_k + \sum_i g_i(k)$.
\end{enumerate}

Since $\sum_i g_i(0) = 0$, the secret is unchanged: $s' = s$. But shares are
fully re-randomized: knowledge of old shares provides no information about new
shares.

Hanzo infrastructure performs key refresh every 24 hours. Refresh completes in
$<$100ms and does not interrupt signing operations.

% ============================================================================
\section{Integration with Hanzo Services}
\label{sec:services}
% ============================================================================

\subsection{IAM Token Signing}

Hanzo IAM~\cite{hanzoiam} issues JWT tokens signed with Ed25519. The IAM
signing key is split across 5 threshold signing nodes in a $(3, 5)$
configuration:

\begin{enumerate}[leftmargin=1.4em]
    \item User authenticates via OIDC (password, passkey, or SSO).
    \item IAM constructs the JWT payload (claims, expiry, org scope).
    \item IAM sends the signing request to the threshold signing service.
    \item Three of five signing nodes produce partial signatures; the
    coordinator aggregates them into a standard Ed25519 signature.
    \item The JWT is returned to the user. Verifiers use the public key
    from the IAM JWKS endpoint---no knowledge of threshold signing needed.
\end{enumerate}

\subsection{KMS Master Key Protection}

Hanzo KMS~\cite{hanzokms} uses envelope encryption: data encryption keys (DEKs)
are encrypted with a master key (KEK). The KEK is threshold-shared:
\begin{itemize}[leftmargin=1.1em]
    \item KEK encryption uses threshold ECDH: each signing node holds a share
    of the ECDH private key, and $t$ nodes must cooperate to derive the shared
    secret used for AES-256-GCM encryption of the DEK.
    \item KEK rotation is a key refresh operation---existing DEKs encrypted
    under the old KEK are re-encrypted under the refreshed key without
    exposing plaintext DEKs.
\end{itemize}

\subsection{Agent Credential Signing}

AI agents~\cite{hanzoagentsdk} receive delegated credentials via certificate
chains:
\begin{enumerate}[leftmargin=1.4em]
    \item The platform root key (threshold-shared) signs an org-level
    intermediate certificate.
    \item The org intermediate signs an agent-specific credential with
    constrained permissions (allowed tools, API scopes, time bound).
    \item Agents present the credential chain when accessing resources. Each
    service verifies the chain back to the well-known root public key.
\end{enumerate}

\subsection{ACI Compute Attestation}

ACI GPU nodes~\cite{hanzoaci} sign inference attestations with threshold ECDSA:
\begin{itemize}[leftmargin=1.1em]
    \item Each GPU node in a verification group holds a share of the group
    signing key.
    \item After verifying an inference result, $t$ of $n$ verifiers produce a
    threshold ECDSA signature over the result hash.
    \item The signature is submitted on-chain as proof of verified computation.
\end{itemize}

% ============================================================================
\section{Security Analysis}
\label{sec:security}
% ============================================================================

\begin{theorem}[Threshold Security]
The Hanzo Threshold Signing system with $(t, n) = (3, 5)$ provides security
under the following adversary model: an adversary who compromises up to $t-1 = 2$
signing nodes cannot forge signatures, recover the private key, or distinguish
partial signatures from random values.
\end{theorem}

\begin{remark}
The identifiable abort property of CGGMP21 ensures that a compromised node
cannot cause signing to fail silently. If a node submits invalid partial
signatures, the coordinator detects the fault, identifies the malicious node,
and re-runs the protocol with the remaining honest nodes (provided $\geq t$
honest nodes remain).
\end{remark}

\subsection{Deployment Topology}

The 5 signing nodes are deployed across 3 geographic regions and 2 cloud
providers:
\begin{itemize}[leftmargin=1.1em]
    \item 2 nodes in US (DOKS lux-k8s, separate availability zones).
    \item 2 nodes in EU (DOKS lux-k8s, separate availability zones).
    \item 1 node in APAC (DOKS lux-k8s).
\end{itemize}

An adversary must compromise nodes in at least 2 geographic regions to reach
the threshold.

% ============================================================================
\section{Performance}
\label{sec:performance}
% ============================================================================

\begin{table}[H]
\centering
\small
\begin{tabular}{lrrr}
\toprule
\textbf{Operation} & \textbf{Protocol} & \textbf{Median (ms)} & \textbf{p99 (ms)} \\
\midrule
DKG (one-time) & FROST & 85 & 120 \\
DKG (one-time) & CGGMP21 & 2400 & 3100 \\
Signing & FROST & 12 & 28 \\
Signing (pre-computed nonce) & FROST & 8 & 18 \\
Signing & CGGMP21 & 45 & 95 \\
Key refresh & FROST & 35 & 60 \\
Key refresh & CGGMP21 & 180 & 310 \\
\bottomrule
\end{tabular}
\caption{Latency for $(3, 5)$ threshold operations across 3 geographic regions.
Measured over 30 days of production traffic.}
\label{tab:perf}
\end{table}

The system processes over 2 million threshold signatures per day (predominantly
FROST for IAM tokens), with 99.97\% availability over the measurement period.

% ============================================================================
\section{Related Work}
\label{sec:related}
% ============================================================================

\textbf{FROST}~\cite{frost} provides the theoretical foundation for our Ed25519
threshold scheme. We extend it with pre-computed nonce batching for
latency-sensitive IAM signing.

\textbf{CGGMP21}~\cite{cggmp21} improves upon GG20~\cite{gg20} with
identifiable abort and reduced communication complexity. We adopt CGGMP21 for
all ECDSA operations.

\textbf{Lux Network}~\cite{luxwhitepaper} uses threshold signatures for
validator set management and cross-chain bridge security. Hanzo Threshold
Signing reuses the same FROST implementation for infrastructure signing.

\textbf{Fireblocks MPC}~\cite{fireblocks} provides commercial MPC wallet
infrastructure. Unlike Fireblocks, Hanzo Threshold Signing is purpose-built for
AI infrastructure signing (IAM, KMS, agent credentials), not financial custody.

% ============================================================================
\section{Conclusion}
\label{sec:conclusion}
% ============================================================================

Hanzo Threshold Signing eliminates single points of compromise across Hanzo
infrastructure by distributing signing keys via $(3, 5)$ threshold schemes.
FROST for Ed25519 achieves 8--12ms signing latency for IAM tokens and agent
credentials. CGGMP21 for ECDSA achieves 45ms signing latency for blockchain
operations. Proactive key refresh every 24 hours limits exposure from share
compromise. The system processes 2M+ daily threshold signatures with 99.97\%
availability.

% ============================================================================
% BIBLIOGRAPHY
% ============================================================================
\begin{thebibliography}{15}

\bibitem{hanzoiam}
{Hanzo AI Research}.
\newblock Hanzo {IAM}: Identity and access management for {AI} platforms.
\newblock \emph{Hanzo AI Research}, 2024.

\bibitem{hanzokms}
{Hanzo AI Research}.
\newblock Hanzo {KMS}: Secrets management with per-environment rotation and audit trail.
\newblock \emph{Hanzo AI Research}, September 2021.

\bibitem{hanzoagentsdk}
{Hanzo AI Research}.
\newblock Hanzo Agent {SDK}: Framework for building {AI} agents.
\newblock \emph{Hanzo AI Research}, 2025.

\bibitem{hanzoaci}
{Hanzo AI Research}.
\newblock Hanzo {ACI}: {AI} Chain Infrastructure for decentralized compute markets.
\newblock \emph{Hanzo AI Research}, December 2024.

\bibitem{luxwhitepaper}
{Lux Partners}.
\newblock Lux: A scalable multi-consensus blockchain with post-quantum cryptography.
\newblock \emph{Lux Network Technical Report}, 2024.

\bibitem{frost}
C.~Komlo and I.~Goldberg.
\newblock {FROST}: Flexible round-optimized {Schnorr} threshold signatures.
\newblock In \emph{Proc.\ SAC}, 2020.

\bibitem{cggmp21}
R.~Canetti, R.~Gennaro, S.~Goldfeder, N.~Makriyannis, and U.~Peled.
\newblock {UC} non-interactive, proactive, threshold {ECDSA} with identifiable aborts.
\newblock In \emph{Proc.\ CCS}, 2020.

\bibitem{pedersen1991}
T.~Pedersen.
\newblock Non-interactive and information-theoretic secure verifiable secret sharing.
\newblock In \emph{Proc.\ CRYPTO}, 1991.

\bibitem{herzberg1995}
A.~Herzberg, S.~Jarecki, H.~Krawczyk, and M.~Yung.
\newblock Proactive secret sharing, or: How to cope with perpetual leakage.
\newblock In \emph{Proc.\ CRYPTO}, 1995.

\bibitem{gg20}
R.~Gennaro and S.~Goldfeder.
\newblock One round threshold {ECDSA} with identifiable abort.
\newblock \emph{IACR Cryptology ePrint Archive}, 2020/540.

\bibitem{fireblocks}
{Fireblocks}.
\newblock {MPC} wallet infrastructure for digital assets.
\newblock \emph{Fireblocks Documentation}, 2023.

\end{thebibliography}

\end{document}
